%%%%%%%%%%%%%%%%%%%%%%%%%%%%%%%%%%%%%%%%%%%%%%%%%%%%%%%%%%%%%%%%%%%%%%
% Lớp 11 | Dãy số. Cấp số cộng. Cấp số nhân | Cấp số nhân
% Dạng 1D2?3-3: Số hạng tổng quát của cấp số nhân
% Số câu: 28  (sinh tự động bởi tach_cau.py; từ nay thầy sửa trực tiếp file này)
%%%%%%%%%%%%%%%%%%%%%%%%%%%%%%%%%%%%%%%%%%%%%%%%%%%%%%%%%%%%%%%%%%%%%%

%%%%%%%%%%%%%%%% TRẮC NGHIỆM 4 PHƯƠNG ÁN %%%%%%%%%%%%%%%%

% nguồn: Ok_Toan11_TracNghiem_Lan2
\begin{ex}%[1D2N3-3]%[ma:6a8d99ce]
	Cho cấp số nhân $(u_n)$ với $u_1=2$ và công bội $q=3$. Giá trị của $u_2$ bằng 
	\choice 
		{\True  $6$}
		 {$9$}
		 {$8$}
		 {$\dfrac{2}{3}$}
	\loigiai{ 
	 Ta có $u_2=u_1 \cdot q=2 \cdot 3=6$.}
\end{ex}

% nguồn: Ok_Toan11_TracNghiem_Lan2
\begin{ex}%[1D2N3-3]%[ma:62eb5d39]
	Cho cấp số nhân $(u_n)$ với $u_1=3$ và công bội $q=2$. Số hạng tổng quát  $u_n$ $(n\ge 2)$ bằng
	\choice 
		 {$3\cdot 2^n$}
		{$3\cdot  2^{n+2}$}
		{$3\cdot  2^{n+1}$}
		 {\True $3\cdot 2^{n-1}$}
	 \loigiai{
	Ta  có $u_n=u_1 \cdot q^{n-1}=3 \cdot 2^{n-1}$.}
\end{ex}

% nguồn: Ok_Toan11_TracNghiem_Lan2
\begin{ex}%[1D2N3-3]%[ma:7f845bf5]
	Cho cấp số nhân $(u_n)$ với $u_1=3$  và công bội $q=2$. Giá trị của $u_2$ bằng
	 \choice
		{$8$} 
		 {$9$}
		 {\True $6$}
		{$\dfrac{3}{2}$} 
	 \loigiai{
	Ta có $u_2=u_1 \cdot q=3 \cdot 2=6$.} 
\end{ex}

% nguồn: Ok_TracNghiemLop11_Lan1
\begin{ex}%[1D2N3-3]%[ma:c04ab661]
	Cho cấp  số nhân $(u_n)$ với $u_1=3$ và công bội $q=2$. Công thức số hạng tổng quát của $(u_n)$ là
	\choice 
		 {\True $u_n=3\cdot 2^{n-1}$}
		{$u_n=3\cdot  2^{n+1}$}
		{$u_n=3\cdot 2^{n+2}$} 
		{$u_n=3\cdot 2^n$} 
	 \loigiai{
		Công thức số hạng tổng  quát của $(u_n)$ là $u_n=u_1\cdot q^{n-1}=3\cdot 2^{n-1}$.
	} 
\end{ex}

% nguồn: Ok_TracNghiemLop11_Lan1
\begin{ex}%[1D2N3-3]%[ma:bb0c8716]
	Cho cấp số nhân $\left(u_n\right)$ có số hạng đầu là $u_1 \neq 0$ và công bội $q \neq 0$. Số hạng tổng quát của cấp số  nhân bằng
	 \choice
		 {$u_n = u_1 + \left(n-1\right)q$}
		{\True $u_n =  u_1 \cdot q^{n-1}$}
		 {$u_n = u_1 \cdot q^{n}$}
		{$u_n =  u_1 \cdot q^{n+1}$}
	 \loigiai{
		Ta có số hạng tổng quát của cấp số nhân là $u_n = u_1  \cdot q^{n-1}$.
	} 
\end{ex}

% nguồn: Ok_TracNghiemLop11_Lan1
\begin{ex}%[1D2N3-3]%[ma:8708a6d9]
	Cho cấp số nhân $(u_n)$ biết $u_2=2$ và $u_3=5$. Khi đó $u_1$  và $q$ lần lượt bằng
	 \choice
		 {$\dfrac{5}{4}$và $\dfrac{5}{2}$}
		{$\dfrac{2}{5}$ và $\dfrac{5}{2}$} 
		 {\True $\dfrac{4}{5}$ và $\dfrac{5}{2}$}
		 {$\dfrac{4}{5}$ và $2$}
	\loigiai{ 
		Ta có $u_3=u_2\cdot q\Rightarrow  q=\dfrac{u_3}{u_2}=\dfrac{5}{2}$.\\
	$u_2=u_1.q\Rightarrow u_1=\dfrac{u_2}{q}=\dfrac{4}{5}$.} 
\end{ex}

% nguồn: Ok_Toan11_TracNghiem_Lan2
\begin{ex}%[1D2H3-3]%[ma:de87176f]
	Trong các dãy số $(u_n)$ cho bởi số hạng tổng  quát $u_n$ sau, dãy số nào là một cấp số nhân?
	 \choice
		{$u_n=7-3n$} 
		{$u_n=7-3^n$} 
		{$u_n=\dfrac{7}{3n}$} 
		{\True $u_n=7\cdot 3^n$} 
	 \loigiai{
		Dãy $u_n=7\cdot 3^n$ là  cấp số nhân có $\heva{&u_1=21 \\&q=3}$.
	 }
\end{ex}

% nguồn: Ok_Toan11_TracNghiem_Lan2
\begin{ex}%[1D2H3-3]%[ma:c936a5e5]
	Trong các dãy số $(u_n)$  có công thức của số hạng tổng quát $u_n$ sau đây, dãy số nào là cấp số nhân?
	\choice 
		 {$u_n=7-3n$}
		 {$u_n=7-3^n$}
		{$u_n=\dfrac{7}{3n}$} 
		{\True  $u_n=7\cdot 3^n$}
	\loigiai{ 
		Dãy số  $(u_n)$ với $u_n=7\cdot 3^n$ có $\dfrac{u_{n+1}}{u_n}=\dfrac{7\cdot 3^{n+1}}{7\cdot 3^n}=3^1=3$.\\
		Do đó $u_n=7\cdot  3^n$ là cấp số nhân với công bội $q=3$ và $u_1=21$.
	 }
\end{ex}

% nguồn: Ok_Toan11_TracNghiem_Lan2
\begin{ex}%[1D2H3-3]%[ma:5fc4a294]
	 Cho dãy số $\left(u_n\right)$ biết $\heva{&u_1=3\\&u_{n+1}=3u_n}, \forall n\in{N^*}$. Tìm số hạng tổng quát của dãy số $(u_n)$.
	 \choice
		{\True $u_n=3^n$} 
		 {$u_n=n^{n+1}$}
		 {$u_n=3^{n+1}$}
		 {$u_n=3^{n-1}$}
	 \loigiai{
		Ta có  $u_1=3$ và $\dfrac{u_{n+1}}{u_n}=3$\\
		Suy ra dãy số $\left(u_n\right)$ là cấp số nhân với  $\heva{&u_1=3\\ &q=3.}$\\
	Do đó $u_n=u_1 \cdot  q^{n-1}=3 \cdot 3^{n-1}=3^n$.}
\end{ex}

% nguồn: Ok_Toan11_TracNghiem_Lan2
\begin{ex}%[1D2H3-3]%[ma:320f4fd2]
	Cho dãy số có các số hạng đầu là $\dfrac{1}{3}$; $\dfrac{1}{3^2}$; $\dfrac{1}{3^3}$; $\dfrac{1}{3^4}$; $\ldots$ Số hạng tổng quát của dãy số  này là
	 \choice
		{$\dfrac{1}{3^{n-1}}$} 
		 {$\dfrac{1}{3^{n+2}}$}
		{\True  $\dfrac{1}{3^n}$}
		{$\dfrac{1}{3^{n+1}}$} 
	\loigiai{ 
		Ta  có\\
		${u_1}=\dfrac{1}{3}=\dfrac{1}{3^1}\\ 
		{u_2}=\dfrac{1}{3^2}=\dfrac{1}{3^2}\\ 
		{u_3}=\dfrac{1}{3^3}=\dfrac{1}{3^3}\\ 
		 \ldots$\\
	Vậy $u_n=\dfrac{1}{3^n}$.} 
\end{ex}

% nguồn: Ok_Toan11_TracNghiem_Lan2
\begin{ex}%[1D2H3-3]%[ma:8719be74]
	Cho một cấp số nhân có các số hạng đều không âm thỏa mãn $u_2=6$, $u_4=24$. Tính tổng của $12$ số hạng đầu  tiên của cấp số nhân đó.
	 \choice
		{\True $3\cdot  2^{12}-3$}
		 {$2^{12}-1$}
		{$3\cdot  2^{12}-1$}
		{$3\cdot  2^{12}$}
	\loigiai{ 
		Gọi công bội của cấp số nhân  bằng $q$.\\
		Suy ra $u_4=u_2\cdot q^2$ $\Rightarrow q=\pm  2$.\\
		Do CSN có các số hạng không âm  nên $q=2$.\\
		Ta có $S_{12}=u_1\cdot \dfrac{1-q^{12}}{1-q}$ $=3\cdot  \dfrac{1-2^{12}}{1-2}=3\left(2^{12}-1\right)$.
	 }
\end{ex}

% nguồn: Ok_TracNghiemLop11_Lan1
\begin{ex}%[1D2H3-3]%[ma:3133d4a0]
	Cho cấp số nhân $3, -\dfrac{3}{4}, \dfrac{3}{16}$. Khi đó  số hạng $u_{n}$ của cấp số nhân bằng
	\choice 
		{$u_n = \left(-\dfrac{1}{4}\right)^{n  + 1} $}
		{$u_n =  \left(-\dfrac{1}{4}\right)^{n - 1} $}
		{\True $u_n = 3 \cdot \left(-\dfrac{1}{4}\right)^{n - 1}  $}
		{$u_n  = 3 \cdot \left(-\dfrac{1}{4}\right)^n$}
	 \loigiai{
	Ta có $3, -\dfrac{3}{4}, \dfrac{3}{16}$ suy ra số hạng đầu $u_1 = 3$ và công bội bằng $q = -\dfrac{1}{4}$. Do đó $u_n = 3 \cdot \left(-\dfrac{1}{4}\right)^{n - 1}  $.}
\end{ex}

% nguồn: Ok_TracNghiemLop11_Lan1
\begin{ex}%[1D2H3-3]%[ma:a3e8b6b3]
	Cho cấp số nhân $(u_n)$ với $u_1=1$; $q=\dfrac{1}{10}$. Hỏi $\dfrac{1}{10^{103}}$ là số hạng  thứ bao nhiêu?
	\choice 
		 {$105$}
		{$103$} 
		 {\True $104$}
		{$102$} 
	\loigiai{ 
		Ta có số hạng thứ $n$ là $u_n=u_1 \cdot q^{n-1}=1 \cdot \left(\dfrac{1}{10}\right)^{n-1}=\dfrac{1}{10^{n-1}}$.\\ 
		Mà $u_n=\dfrac{1}{10^{103}}$  suy ra $n-1=103$ hay $n=104$.
	} 
\end{ex}

% nguồn: Ok_Toan11_TracNghiem_Lan2
\begin{ex}%[1D2V3-3]%[ma:3d286ca3]
	Cho dãy số $\left(u_n\right)$ thỏa  mãn $\heva{&3u_1=\sqrt{3u_1-u_2}+u_2+6\\ &u_{n+1}=2u_n},\forall n\in {\mathbf{N}^*}$. Tìm giá trị nhỏ nhất của $n$ để $u_n\ge{2^{2021}}$.
	 \choice
		 {$2021$}
		 {$1012$}
		 {$2022$}
		 {\True $1011$}
	\loigiai{ 
		Ta có $u_{n+1}=2u_n \Rightarrow \dfrac{u_{n+1}}{u_n}=2,\forall n\in{\mathbf{N}^*}$ nên dãy  $(u_n)$ là cấp số nhân với công bội $q=2$.\\
		 $\Rightarrow{u_2}=2u_1$.\\
		Mà 
		\begin{eqnarray*}
			&&3u_1=\sqrt{3u_1-u_2}+u_2+6\\ 
			&\Leftrightarrow& 3u_1-u_2-\sqrt{3u_1-u_2}-6=0\\ 
			&\Leftrightarrow&{\left(\sqrt{3u_1-u_2}\right)^2}-\sqrt{3u_1-u_2}-6=0\\ 
			&\Leftrightarrow& \hoac{&\sqrt{3u_1-u_2}=2 \qquad \text{(Nhận)}\\ &\sqrt{3u_1-u_2}=-3  \qquad \text{(Loại)}}\\
			&\Leftrightarrow& 3u_1-u_2=4. 
		\end{eqnarray*}
		Từ và ta có $\heva{ &{u_2}=2u_1\\ & 3u_1-u_2=4}  \Rightarrow{u_1}=4$.\\
		$\Rightarrow\left(u_n\right)$ là cấp số nhân  với công bội $q=2$, $u_1=4$.\\
		Nên số hạng tổng quát  là $u_n=2 \cdot 4^{n-1}=2 \cdot 2^{2(n-1)}=2^{2n-1},\forall n\in{\mathbf{N}^*}$.\\
		$u_n\ge{2^{2021}}\Leftrightarrow{2^{2n-1}}\ge{2^{2021}} \Leftrightarrow 2n-1\ge  2021\Leftrightarrow n\ge 1011$.\\
	Vậy giá trị nhỏ nhất thỏa  mãn là $1011$.}
\end{ex}

% nguồn: Ok_Toan11_TracNghiem_Lan2
\begin{ex}%[1D2V3-3]%[ma:f36aa6b9]
	Cho dãy  số $(u_n)$ biết $\heva{&u_1=3\\&u_{n+1}=3u_n},\forall n\in \mathbb{N}^*$. Tìm số hạng tổng quát của dãy số $(u_n)$.
	\choice 
		{$u_n=n^{n+1}$} 
		{\True $u_n=3^n$} 
		{$u_n=3^{n+1}$} 
		 {$u_n=3^{n-1}$}
	\loigiai{ 
		Ta có $(u_n)$ là cấp số nhân với số hạng đầu $u_1=3$  và công bội $q=\dfrac{u_{n+1}}{u_n}=3$.\\
	 Do đó số hạng tổng quát $u_n=u_1\cdot q^{n-1}=3\cdot 3^{n-1}=3^n$.}
\end{ex}

% nguồn: Ok_Toan11_TracNghiem_Lan2
\begin{ex}%[1D2V3-3]%[ma:bbd2df55]
	Cho dãy số $\left(u_n \right)$ thỏa  mãn
	$\heva{&5u_1+\sqrt{5u_1-u_2}=u_2+6 \\ &u_{n+1}=3u_n}$,  $\forall n\in \mathbb{N^*}$.
	Giá trị nhỏ nhất của $n$ để $u_n  \ge 2\cdot 3^{2\,018}$ bằng
	\choice 
		{$2\,017$} 
		 {$2\,018$}
		{\True $2\,019$} 
		 {$2\,010$}
	\loigiai{ 
		 \[\heva{&{5u_1+\sqrt{5u_1-u_2}=u_2+6\qquad(1)} \\
		 &{u_{n+1}=3u_n ,\, \forall n\in \mathbb{N}^{*}\quad\qquad\qquad (2)}
		.}\] 
		Từ $(1)$ có 
		\begin{align*}
			5u_1+\sqrt{5u_1-u_2}=u_2+6 
			&\Leftrightarrow \left(\sqrt{5u_1-u_2} \right)^2+\sqrt{5u_1-u_2}-6=0\\ 
			&\Leftrightarrow  \sqrt{5u_1-u_2}=2\Leftrightarrow 5u_1-u_2=4.
		\end{align*}
		Từ  $(2)$ có $u_{n+1}=3u_n \Rightarrow u_2=3u_1$.\\
		 Giải hệ $\heva{&{5u_1-u_2=4} \\
		&{u_2=3u_1}}$ được $u_1=2$.\\ 
		Dãy $\left(u_n \right)$ là cấp số  nhân với $\heva{&{u_1=2} \\
		&{q=3}}$  có số hạng tổng quát $u_n=2\cdot 3^{n-1}$ với $n\in \mathbb{N}^*$.\\
		\[u_n \ge  2\cdot 3^{2\,018} \Leftrightarrow 2\cdot 3^{n-1} \ge 2\cdot 3^{2\,018} \Leftrightarrow n-1\ge 2\,018\Leftrightarrow n\ge 2\,019.\]
		Vậy giá trị nhỏ nhất thỏa mãn là $2\,019$. 
	} 
\end{ex}

% nguồn: Ok_Toan11_TracNghiem_Lan2
\begin{ex}%[1D2V3-3]%[ma:6334365b]
	 Cho dãy số $(u_n)$ thỏa mãn $\heva{&u_1=1\\&u_n=2u_{n-1}+1;~n\ge 2}$. Tổng $S=u_1+u_2+\cdots+u_{20}$ bằng
	 \choice
		{$2^{20}-20$} 
		{\True  $2^{21}-22$}
		 {$2^{20}$}
		{$2^{21}-20$} 
	\loigiai{ 
		Đặt $v_n=u_n+1$, ta có $v_n=2v_{n-1}$ trong  đó $v_1=2$.\\
		Vậy $(v_n)$  là cấp số nhân có số hạng đầu $v_1=2$ và công bội bằng $2$, nên số hạng tổng quát
		\[v_n=2^n \Rightarrow u_n=v_n-1=2^n-1.\] 
		\begin{eqnarray*}
			\Rightarrow S &=&  u_1+u_2+\cdots+u_{20}\\
			 &=& \left(2^1-1\right)+\left(2^2-1\right)+\cdots+\left(2^{20}-1\right)\\
			&=& \left(2^1+2^2+\cdots+2^{20}  \right)-20\\
			&=&  2\left(2^{20}-1\right)-20\\
			&=& 2^{21}-22. 
		\end{eqnarray*}
	 }
\end{ex}

% nguồn: Ok_Toan11_TracNghiem_Lan2
\begin{ex}%[1D2V3-3]%[ma:ad216bcd]
	 Cho dãy số có các số hạng đầu là $\dfrac{1}{3}, \dfrac{1}{3^2}, \dfrac{1}{3^3}, \dfrac{1}{3^4}, \ldots$. Số hạng tổng quát của dãy số này là
	\choice 
		{ $\dfrac{1}{3^{n-1}}$ } 
		{ $\dfrac{1}{3^{n+2}}$ } 
		{\True  $\dfrac{1}{3^n}$ }
		 { $\dfrac{1}{3^{n+1}}$ }
	\loigiai{ 
		Ta có: $u_1 = \dfrac{1}{3^1}$, $u_2 = \dfrac{1}{3^2}$,  $u_3 = \dfrac{1}{3^3}$, $\ldots$ \\
		Vậy $u_n  = \dfrac{1}{3^n}$.
	 }
\end{ex}


%%%%%%%%%%%%%%%% ĐÚNG/SAI %%%%%%%%%%%%%%%%

% nguồn: Ok_dungsai_lop11_Lan1
\begin{ex}%[1D2B3-3]%[ma:92341d7c]
  Hãy cho biết tính đúng sai của mỗi mệnh đề dưới đây?
 \choiceTF 
  {\True Cho cấp số  nhân $\left(u_n\right)$ với $u_1=1$, $q=2$. Giá trị $u_4$ bằng $8$}
  {Cho cấp số  nhân $(u_n)$ có $u_1=-\dfrac{1}{2}$ và $u_2=1$. Giá trị $u_{10}$ bằng $512$}
  {\True Cho cấp số nhân $\left(u_n \right)$, biết $u_1=14$ và $u_2=-28$.  Giá trị $u_{3}$ bằng $56$}
  {Cấp số nhân $(u_n)$ có $u_1=-3$, $\dfrac{u_8}{u_5}=125$.  Giá trị $u_3$ bằng $-375$}
 \loigiai{ 
  \begin{itemchoice}
   \itemch Ta  có $u_4=u_1q^3=1\cdot 2^3=8$.
   \itemch Cấp số nhân có công bội $q=\dfrac{u_2}{u_1}=-2$. Suy ra  $u_{10}=u_1q^9=-\dfrac{1}{2}\cdot(-2)^9=256.$
   \itemch Ta có $u_1 u_3 = u_2^2 \Leftrightarrow u_3 = \dfrac{u_2^2}{u_1}=\dfrac{(-28)^2}{14}=56$. 
   \itemch Ta có $\dfrac{u_8}{u_5}=125\Leftrightarrow \dfrac{u_1\cdot  q^7}{ u_1\cdot q^4}=125\Leftrightarrow q^3=125\Leftrightarrow q=5$.\\
   Vậy $u_3=u_1\cdot  q^2=-3\cdot 5^2=-75$.
  \end{itemchoice}
  }
\end{ex}

% nguồn: Ok_dungsai_lop11_Lan1
\begin{ex}%[1D2B3-3]%[ma:14432319]
 Hãy cho biết tính đúng sai của mỗi mệnh  đề dưới đây?
 \choiceTF 
  {Cho cấp số nhân $1,3,9,\dots$  Số hạng tổng quát của cấp số nhân đó là $u_n=3^{n+1}$}
  {\True Cho dãy số có các số hạng đầu là  $1$; $5$; $25$; $125$; $625$; $\ldots$ . Số hạng tổng quát của dãy số này là $u_n=5^{n-1}$}
  {Cho dãy số có các số hạng đầu là $\dfrac{1}{3}; \dfrac{1}{3^{3}}; \dfrac{1}{3^{3}}; \dfrac{1}{3^4}; \dfrac{1}{3^5}; \ldots $. Số  hạng tổng quát của dãy số là $u_{n}=\dfrac{1}{3^{n-1}}$}
  {\True  Cho dãy số $(u_{n})$ thỏa mãn $\heva{&u_{1}=10\\ &u_{n+1}=2u_{n}}$ ($n\in\mathbb{N}^{*}$). Số hạng tổng quát của dãy số là $u_{n}=10\cdot 2^{n-1}$}
  \loigiai{
  \begin{itemchoice}
   \itemch Nhận thấy dãy $1,3,9,\dots$ là một cấp số nhân với $u_1=1$ công bội $q=3$.\\ 
   Do đó số hạng  tổng quát của dãy là $u_n=u_1q^{n-1}=3^{n-1}$.
   \itemch Ta thấy dãy số đã cho là một cấp số nhân với  công bội $q=5$ và $u_1=1$ nên số hạng tổng quát của dãy số này là $u_n=u_1\cdot q^{n-1}=5^{n-1}$.
   \itemch Công thức số hạng tổng quát của dãy trên  là $u_{n}=\dfrac{1}{3^{n}}$.
   \itemch Ta có $u_{n+1}=2u_{n}\Leftrightarrow \dfrac{u_{n+1}}{u_{n}}=2$, suy ra dãy số $(u_{n})$ là một cấp số nhân với $u_{1}=10$ và công bội  $q=2$.\\
    Do đó số hạng tổng quát là $u_{n}=10\cdot 2^{n-1}$.
  \end{itemchoice}
 } 
\end{ex}


%%%%%%%%%%%%%%%% TRẢ LỜI NGẮN %%%%%%%%%%%%%%%%

% nguồn: Ok_TraLoiNgan_Lop11_Lan2
\begin{ex}%[1D2H3-3]%[ma:2e5aaecd]
	Cho cấp số nhân $(u_n)$  có công bội $q$ và thỏa mãn $\heva{&u_1 + u_2 + u_3 = 26\\ &u_1^2 + u_2^2 + u_3^2 = 364.}$\\
	Tìm  $q$ biết rằng $q<-1$.
	\shortans{-3} 
	 \loigiai{
		Ta  có
		 \[
		\heva{&u_1 + u_2 + u_3 = 26\\ &u_1^2  + u_2^2 + u_3^2 = 364}
		 \Leftrightarrow
		\heva{&u_1(1 + q + q^2) = 26\\ &u_1^2(1 + q^2 + q^4) =  364}
		 \Rightarrow
		\heva{&u_1^2(1 + q + q^2)^2 = 26^2 \quad (1)\\ &u_1^2(1 + q^2 + q^4) = 364.  \quad (2)}
		\] 
		Lấy $(1)$ chia  $(2)$, ta được:
		\[ 
		\dfrac{(1 + q + q^2)^2}{1 + q^2 +  q^4} = \dfrac{26^2}{364}
		\Rightarrow 3q^4 - 7q^3 - 4q^2 - 7q + 3  = 0
		\Leftrightarrow 3\left(q^2 +  \dfrac{1}{q^2}\right) - 7\left(q + \dfrac{1}{q}\right) - 4 = 0.
		\] 
		Đặt $t = q + \dfrac{1}{q}$, $|t| \geq 2$. Phương trình trở thành $3t^2 - 7t - 10 = 0 \Rightarrow  \hoac{&t = -1 \text{ (loại)}\\&t = -\dfrac{10}{3}.}$\\
		Với $t = -\dfrac{10}{3}$, suy ra $q + \dfrac{1}{q} = -\dfrac{10}{3}  \Rightarrow 3q^2 + 10q + 3 = 0 \Leftrightarrow q = -3$ hoặc $q =- \dfrac{1}{3}$.\\
		Vì $q <- 1$ nên $q = -3$. 
	} 
\end{ex}


%%%%%%%%%%%%%%%% TỰ LUẬN %%%%%%%%%%%%%%%%

% nguồn: Ok_TuanLuan_Toan11_Lan2
\begin{ex}%[1D2H3-3]%[ma:277eccd7]
	Cho cấp số  nhân $\left(u_n\right)$ thỏa mãn
	$\heva{&u_4=\dfrac{2}{27}\\& u_3=243 u_8.}$ 
	\begin{enumerate}
		\item Viết năm số hạng đầu  của cấp số nhân.
		\item Số $\dfrac{2}{6561}$  là số hạng thứ bao nhiêu của cấp số nhân?
	\end{enumerate}
	\loigiai{ 
		Gọi $q$ là  công bội của cấp số. Theo giả thiết ta có
		\[\heva{&u_1 q^3=\dfrac {2}{2 7}\\&u_1q^2=243u_1q^7} \Leftrightarrow \heva{&u_1 q^3=\dfrac {2}{2 7}\\&q ^ {5}=\dfrac {1}{2 4 3}} \Leftrightarrow \heva{& q=\dfrac {1}{3}\\&u_1=2.} 
		 \]
		\begin{enumerate}
			\item Năm số hạng đầu của cấp số là: $u_1=2$, $u_2=\dfrac{2}{3}$, $u_3=\dfrac{2}{9}$; $u_4=\dfrac{2}{27}$,  $u_5=\dfrac{2}{81}$.
			\item Ta có $u_n=\dfrac{2}{3^{n-1}} \Rightarrow u_n=\dfrac{2}{6561}  \Leftrightarrow 3^{n-1}=6561=3^8 \Rightarrow n=9$.\\
			Vậy $\dfrac{2}{6561}$  là số hạng thứ $9$ của cấp số nhân.
		\end{enumerate}
	} 
\end{ex}

% nguồn: Ok_TuanLuan_Toan11_Lan2
\begin{ex}%[1D2H3-3]%[ma:18b5ead7]
	Cho cấp số nhân $\left(u_n\right)$ biết $u_1+u_5=51; u_2+u_6=102$. Hỏi số 12288  là số hạng thứ mấy của cấp số nhân $\left(u_n\right)$ ?
	\loigiai{ 
		Gọi $q$ là công bội của cấp số nhân đã cho.  Theo đề bài, ta có
		\[ 
		\left\{\begin{array}{l}{u _ {1}+u _ {5}=5 1} \\
		{u _ {2}+u _ {6}=1  0 2}
	\end{array} \Leftrightarrow \left\{\begin{array}{l}
	u_1\left(1+q^4\right)=51 \\ 
	u_1  q\left(1+q^4\right)=102
	\end{array} \Rightarrow q=2 \Rightarrow u_1=3 \Rightarrow u_n=3.2^{n-1}\right.\right.
	\] 
	Mặt khác $u_n=12288 \Leftrightarrow 3.2^{n-1}=12288 \Leftrightarrow  2^{n-1}=2^{12} \Leftrightarrow n=13$. }
\end{ex}

% nguồn: Ok_TuanLuan_Toan11_Lan2
\begin{ex}%[1D2H3-3]%[ma:e8ecd6dc]
	 Cho $5$ số lập thành một cấp số nhân. Biết công bội bằng một phần số hạng đầu tiên và tổng $2$ số hạng đầu bằng $8$.
	\loigiai{ 
		\[\heva{&u_1+u_2=8\\&q=\dfrac{1}{4}u_1} \Leftrightarrow  \heva{&u_1^2+4u_1=32\\&q=\dfrac{1}{4}u_1} \Leftrightarrow \heva{&\hoac{&u_1=-8\\&u_1=4}\\&\hoac{&q=-2\\&q=1}.}\]
		Vậy CSN là $-8$, $16$, $-32$,  $64$, $-128$ hoặc $4$, $4$, $4$, $4$, $4$.
	 }
\end{ex}

% nguồn: Ok_TuanLuan_Toan11_Lan2
\begin{ex}%[1D2H3-3]%[ma:3a2c9414]
	Cho tứ giác $ABCD$ có $4$ góc tạo thành  $1$ cấp số nhân có công bội bằng $2$. Tìm $4$ góc ấy
	\loigiai{ 
		\[\heva{&u_1+u_2+u_3+u_4=360^\circ\\&q=2} \Leftrightarrow  \heva{&u_1\dfrac{1-q^4}{1-q}=360^\circ\\&q=2} \Leftrightarrow \heva{&u_1=42^\circ\\&q=2.}\]
		 Vậy 4 góc là $24^\circ$, $48^\circ$, $96^\circ$, $192^\circ$.
	} 
\end{ex}

% nguồn: Ok_TuanLuan_Toan11_Lan2
\begin{ex}%[1D2H3-3]%[ma:2a200bd4]
	Cho cấp số nhân  $\left(u_n\right)$ với công bội $q < 0$ và $u_2=4, u_4=9$. Tìm $u_1$.
	 \loigiai{
		Vì $q < 0, u_2 >  0$ nên $u_3 < 0$. Do đó $u_3=-\sqrt{u_2 \cdot u_4}=-\sqrt{4.9}=-6$;
		$u_2^2=u_1 \cdot  u_3 \Rightarrow u_1=\dfrac{u_2^2}{u_3}=\dfrac{4^2}{-6}=-\dfrac{8}{3}$.
	 }
\end{ex}

% nguồn: Ok_TuanLuan_Toan11_Lan2
\begin{ex}%[1D2V3-3]%[ma:9f5eddd5]
	Viết năm số hạng đầu của mỗi dãy số $(u_n)$ sau và xem nó có phải là cấp số nhân không. Nếu nó là cấp số nhân, hãy tìm công bội $q$ và viết công thức số  hạng tổng quát của nó dưới dạng $u_n = u_1 \cdot q^{n-1}$.
	\begin{listEX}[2]
		\item  $u_n = 5n$
		\item $u_n =  5^n$
		\item $u_1 = 1$, $u_n =  n \cdot u_{n-1}$
		\item $u_1 =  1$, $u_n = 5 \cdot u_{n-1}$
	\end{listEX}
	 \loigiai{
		\begin{listEX}
			\item Năm số hạng đầu của dãy:  $5$; $10$; $15$; $20$; $25$.\\
			Ta có $10 \div 5 = 2 \ne  15 \div 10 = \dfrac{3}{2}$ suy ra $(u_n)$ không phải cấp số nhân.
			\item Năm số hạng  đầu của dãy: $5$; $25$; $125$; $625$; $3125$.\\
			Ta có $u_n = 5^n$ nên $u_{n+1} = 5^{n+1} \Rightarrow \dfrac{u_{n+1}}{u_n} = \dfrac{5^{n+1}}{5^n} = 5,~\forall n \ge  1$.\\
			Do đó $(u_n)$ là cấp số nhân có công bội $q = 5$. Số hạng tổng quát $u_n = 5  \cdot 5^{n-1}$.
			 \item
			Năm  số hạng đầu của dãy: $1$; $2$; $6$; $24$; $120$.\\
			Ta có: $2 \div 1 = 2 \ne  6 \div 2 = 3$ nên $(u_n)$ không phải là cấp số nhân.
			\item  Năm số hạng đầu của dãy $1$; $5$; $25$; $125$; $625$.\\
			Ta có $u_n = 5u_{n-1} \Rightarrow \dfrac{u_n}{u_{n-1}} = 5$; $(\forall n \ge 2)$.\\ 
			Do đó $(u_n)$ là cấp số nhân với công sai $q  = 5$. Số hạng tổng quát $u_n = 5^{n-1}$.
		\end{listEX}
	 }
\end{ex}

% nguồn: Ok_TuanLuan_Toan11_Lan2
\begin{ex}%[1D2V3-3]%[ma:61523c48]
	Cho  dãy số $\left(u_n\right)$ xác định bởi $\heva{&u_1=2\\&u_{n+1}=4u_n+9,~\left(n\geq 1\right).}$
	\begin{listEX}
		\item 
		 Chứng minh dãy số $\left(v_n\right)$ với $v_n=u_n+3$, $n\leq 1$ là một cấp số nhân.
		 \item
		Tìm công thức tổng quát  của dãy số $\left(u_n\right)$.
	\end{listEX}
	\loigiai{ 
		\begin{listEX}
			\item 
			Ta có  $v_n=u_n+3 \Rightarrow v_{n+1}=u_{n+1}+3=\left(4u_n+9\right)+3$ nên
			\begin{eqnarray*}
				 \dfrac{v_{n+1}}{v_n}&=&\dfrac{\left(4u_n+9\right)+3}{u_n}\\
				&=&\dfrac{4\left(u_n+3\right)}{u_n+3}\\ 
				&=&4. 
			\end{eqnarray*}
			Vậy $\left(v_n\right)$ là cấp số nhân với số hạng đầu $v_1=v_1+3=2+3=5$ và công  bội $q=4$.
			 \item
			Do $\left(v_n\right)$ là cấp số nhân với $\heva{&v_1=5\\&q=4}$ nên số hạng tổng quát của $v_n=v_1\cdot q^{n-1}=5\cdot 4^{n-1}$.\\ 
			Suy ra công thức tổng quát của dãy số $\left(u_n\right)$ là  $u_n=v_n-3=5\cdot 4^{n-1}-3$.
		\end{listEX}
	} 
\end{ex}
